% ---------------------------------------------------------------------------
% Benchmark construction. Number-independent: everything here is a property of
% the artifact, not of any judge's behaviour, so it is stable across runs.
% ---------------------------------------------------------------------------
\section{The JudgeSense benchmark}
\label{sec:benchmark}

JudgeSense measures whether an LLM judge returns the same verdict when the same
evaluation request is worded differently. The construct is deliberately narrow:
we are not asking whether a judge is \emph{correct}, but whether it is
\emph{stable}. A judge that is consistently wrong is reproducible; a judge whose
verdict depends on which of two equivalent phrasings it received is not usable as
a measuring instrument, however accurate it looks on average.

\subsection{Tasks and sources}
\label{sec:sources}

The benchmark covers four evaluation tasks, each drawn from an existing
human-labelled corpus. No item is authored by us, and no label is assigned or
inferred by us; every label is the one the source corpus already carried.

\begin{description}
  \item[Factuality] (250 items, TruthfulQA). A statement is judged accurate or
    inaccurate. Items whose question is ill-posed, and whose only defensible
    answer is a refusal to answer, are excluded at load time.
  \item[Coherence] (250 items, SummEval). A summary is rated 1--5 for
    coherence against expert annotations. Summaries with identical text but
    conflicting expert labels are dropped rather than merged.
  \item[Relevance] (250 items, 500 rows, BEIR TREC-COVID). Given a topic, the
    judge picks the more relevant of two passages. \emph{Both} candidates carry
    a human relevance grade: the positive is graded fully relevant (2) and the
    distractor is one a human assessor explicitly rejected (0). Partially
    relevant documents (grade 1) are used for neither role.
  \item[Preference] (130 items, 260 rows, MT-Bench human judgments). The judge
    picks the response humans preferred. Only comparisons with at least two
    annotator votes and a clear majority are used.
\end{description}

The relevance construction is the substantive departure from the usual practice
of drawing distractors from a positives-only collection. There, a negative can
only be a document \emph{absent} from an incomplete relevance judgment set, so a
topically similar distractor may in fact be relevant but unjudged, and the item
penalises a defensible answer. TREC-COVID carries graded judgments including
tens of thousands of explicit non-relevant assessments, so here both candidates
are human-labelled and the distractor is one a human affirmatively rejected.

Every shipped item records its source dataset, split, per-record identifier, and
the source configuration where the source defines one, so any item can be traced
to the row it came from.
Loaders raise rather than substitute when a source is unavailable; there is no
fallback path that could silently generate an item.

\subsection{Paraphrase design}

Each item is issued under two instruction templates that differ in wording and
not in what they ask. Templates are held to a fixed label space: a check run in
continuous integration asserts that no factuality or coherence template alters
the set of admissible answers or inverts their polarity. The pairwise templates
share a two-option answer space by construction and are not covered by that
check.

This constraint is a correction of our own earlier design. An earlier version of
this benchmark included a factuality template that inverted the answer polarity
(``does this response contain factual errors? answer NO if accurate''). Under a
single shared label space, ``YES'' then denoted opposite conclusions in different
arms, so a judge that answered both arms correctly was recorded as having
disagreed with itself. That is a label-convention defect in the benchmark rather
than a property of any model, and it measures a different construct: handling a
polarity flip is a comprehension task, whereas we ask whether equivalent
phrasings yield equivalent verdicts.

We state the direction of the resulting bias rather than leave it implicit.
Polarity-inverted pairings are among the hardest items in the set, so excluding
them can only push measured consistency \emph{upward}; our stability figures are
conditioned on that exclusion and should be read accordingly. The alternative ---
retaining the inverted templates and remapping each template's raw answers onto a
single canonical decision space before comparison --- is implemented in the
released code, and we report the remapped analysis alongside the
excluded-pairing analysis so the exclusion's effect on the reported numbers is
visible rather than assumed.

\subsection{Position control}

Both pairwise tasks present every item twice, once in each candidate ordering.
The correct answer therefore sits at position A in exactly half the rows, and a
judge that ignores content and always answers ``A'' scores $0.500$ rather than
$1.000$.

This too corrects an earlier version, in which the correct candidate always
appeared first. Under that construction a position-anchored judge was
indistinguishable from a perfect one, and most judges scored a degenerate
$\mathrm{JSS}=1.000$ that reflected the layout rather than any judgment.

\subsection{Shortcut controls}
\label{sec:shortcuts}

A benchmark measures what it claims to measure only if it cannot be passed by a
heuristic that ignores the intended construct. We therefore evaluate the shipped
splits against judges that use no judgment at all. Each is run over the released
files, not over the loaders, so these are properties of the artifact:

\begin{itemize}
  \item \textbf{Position only} (always answer A): $0.500$ on relevance
    ($n=500$) and $0.500$ on preference ($n=260$), by construction of the swap
    design.
  \item \textbf{Length only} (pick the longer candidate): $0.482$ on relevance
    ($n=494$), $0.508$ on preference ($n=256$).
  \item \textbf{Lexical overlap only} (pick the candidate sharing more terms with
    the query): $0.540$ on relevance ($n=252$).
\end{itemize}

Support counts fall below the row counts because a heuristic that ties has no
signal to act on; tied rows are excluded rather than split, which is why the
overlap control rests on roughly half the relevance rows. All three controls are
computed from the released files by \texttt{scripts/shortcut\_controls.py}, whose
output is committed, so every figure above is reproducible without rerunning any
judge.

None of these exceeds chance in either direction, though the overlap control is
the weakest of the three and we do not overstate it: its 95\% Wilson interval is
$[0.478, 0.600]$, so the data are consistent with a residual lexical signal worth
about $0.57$. It also reads only half the relevance rows, because the remainder
tie on the overlap statistic and a heuristic with no signal on a row is not
evidence about that row. The length and position controls are far tighter.

The length control is enforced rather than observed. Human annotators in the
preference corpus preferred the longer response in roughly 70\% of usable
comparisons, so an unfiltered sample rewards verbosity directly; we hold the
winner-longer share at exactly 50\% by taking an equal number of items from each
length bucket. That $50\%$ is a construction invariant of the item set and is a
different quantity from the $0.508$ above, which is what a length-following judge
scores on the rendered rows once ties are dropped.

The preference split ships 130 items rather than the round 250, and we state the
full attribution rather than the net gap. Of the source comparisons, 607 failed
the decisive-vote label rule or carried no strict majority, one was duplicate
content, 67 were dropped because their gold contradicted another item's gold on
shared candidate text, and 102 more were dropped to hold the length balance
exactly. Each filter removes a way the benchmark could be passed or contradicted
without judging, and the last two were added after we found both failures in our
own shipped data. We regard the reduced split as the correct trade: the
alternative is a larger set that a verbosity heuristic beats, and one that scores
a consistent judge as wrong whenever the same response appears on both sides of
two different comparisons.

Balancing on length alone would install a second confound in place of the first.
Because the winner-shorter bucket ships whole and only the winner-longer bucket
is cut, any criterion used to make that cut applies to one bucket and not the
other. Cutting by vote margin --- keeping the most contested comparisons, as an
earlier build did --- left the buckets matched on length but far apart on
annotator agreement: mean vote-margin ratio $0.567$ against $0.745$, and $16\%$
against $54\%$ unanimous. A judge that is stronger on close calls would then
read as length-biased. We therefore \emph{stratify}: the winner-shorter bucket
fixes the target distribution of vote shapes (absolute margin and total votes),
and the winner-longer bucket is filled stratum by stratum to match it. The
shipped buckets agree to $0.0003$ in mean vote-margin ratio ($0.8603$ against
$0.8605$), have identical medians, and contain the same number of unanimous
comparisons ($49$ each).

Matching stops the split from collapsing onto coin-flips, but the opposite
failure is equally real and less obvious. The filter that removes contradictory
gold is a greedy selection over a conflict graph, so the order in which items are
considered determines both which survive and how hard the survivors are. Keeping
the firmest human majorities first is the intuitive rule and is worse on both
counts, measured on this pool: $106$ items at $90.6\%$ unanimous, against $130$
items at $75.4\%$ when the most contested are kept first. A split on which the
annotators were nearly always unanimous cannot separate a competent judge from a
mediocre one --- both score near the top --- which is the ceiling effect that
made the earlier version of this benchmark uninformative. We therefore resolve
conflicts contested-first, and report the resulting distribution rather than only
the balance: $75.4\%$ of shipped comparisons carry a unanimous human majority and
the rest are genuine close calls.

The lexical control required a similar correction. Selecting each distractor to
maximise its term overlap with the topic --- the obvious way to make an item
``hard'' --- made distractors \emph{denser} in query terms than the genuinely
relevant passages, so the heuristic ``pick the passage with \emph{lower} overlap''
recovered the answer far above chance. Matching the distractor's retrieval score
to the positive's, rather than maximising it, removes overlap as a signal in
either direction. We report the control two-sided for this reason: a split whose
answer is reliably the \emph{shorter} or \emph{less overlapping} candidate is
exploited exactly as cheaply as one biased the other way.

\subsection{Release and integrity}

The dataset is released with Croissant metadata, per-item provenance, and a
content hash for each split that excludes retrieval timestamps, so a rebuild can
be verified as identical in content without being byte-identical in metadata. A
28-check audit gate runs in continuous integration and fails the build on
duplicate content, unbacked provenance claims, label degeneracy, insufficient
clustering support, contradictory ground truth, or a length shortcut exceeding
its bound. Checks that cannot read the data they are meant to inspect fail
rather than pass, so a check that silently measures nothing cannot be mistaken
for a check that found nothing wrong.
